\section{Planungsphase}
\label{sec:Planungsphase}

\subsection{Zeitlicher Rahmen des Projektes} 
\label{sub:Zeitlicher Rahmen des Projektes} 
Für die Durchführung der betrieblichen Projektarbeit steht ein Zeitrahmen von 70 Arbeitsstunden zur Verfügung. Der
Durchführungszeitraum beginnt am 09. März 2015 und endet am 19. März 2015. Um diesen Zeitraum nicht zu überschreiten ist es
erforderlich, die im Projektantrag festgelegten Meilensteine sowohl zeitlich als auch inhaltlich einzuhalten. Das
Risiko durch abweichende Anforderungen wird mit einem einberechneten Puffer begrenzt.


\subsection{Vorgehensweise} 
\label{sub:Vorgehensweise}
Zur Umsetzung des Projektes wird eine iterativ inkrementelle Vorgehensweise gewählt. Bei diesem
Vorgehen wird das Gesamtprojekt schrittweise entwickelt, also in Teilprojekte (oder auch Iterationen) aufgeteilt. Jedes dieser
Teilprojekte durchlebt dabei alle Entwicklungsdisziplinen, von der Analyse bis zur Implementierung. Nach jeder Iteration kann
die neu gewonnene Erfahrung und möglicherweise erhaltenes Feedback, unmittelbar in die nächste Iteration einfließen.
Darauffolgende Iterationen können auch zur kontinuierlichen Verbesserung vorheriger Teilprojekte dienen.

Diese zyklische Entwicklung soll durch das Einbinden neuster Erkenntnisse, z.B. Anforderungensänderungen, eine hohe
Softwarequalität sichern und die Einhaltung der Anforderungen gewährleisten.


\subsection{Verwendete Hardware} 
\label{sub:Verwendete Hardware}
Neben dem in der \Tabelle{tab:Laptop-Spezifikation} auf S.~\pageref{app:Laptop-Spezifikation} aufgeführten Laptop wurden
eine Dockingstation, mit einer 1:1 Portreplikation, und zwei 24'' Bildschirme (HP Compaq LA2405x) genutzt.


\subsection{Verwendete Software/Technologien} 
\label{sub:Verwendete Software/Technologien}
Die nachfolgende Liste ist aus Gründen der Übersichtlichkeit auf einen Auszug der wichtigsten Spezifikationen reduziert worden.
\begin{description}
	\item [Intellij IDEA] Ist eine kostenpflichtige, dafür sehr mächtige und umfangreiche, integrierte
	Entwicklungsumgebung (von engl.: \Abkuerzung{IDE}).
	\item [Java EE]
	Die Plattform \Abkuerzung{Java EE}, ehemals J2EE ist eine umfangreiche Sammlung von Softwarekomponenten und Diensten. Ferner
	werden Spezifikationen, Code Standards und Bibliotheken definiert. Mit der Java EE Plattform wurde ein Rahmen für
	Softwareentwickler geschaffen, in dem verteilte, transaktions-basierte und mehrschichtige Unternehmensanwendungen entwickelt
	werden können. Klar definierte Bibliotheken und Schnittstellen / \Abkuerzung{API}'s zwischen den Java EE Komponenten sorgen
	dafür, dass auch Software-Komponenten von unterschiedlichen Herstellern zusammenarbeiten können (z.B. SAP Anbindung).
	\citep{JavaEE}
	\item [Atmosphere]
	Das Open-Source-Framework Atmosphere ermöglicht es, in \Abkuerzung{Java EE} skalierbare und asynchrone
	Echtzeit-Anwendungen zu entwickeln. Dabei unterstützt das Framework eine Reihe von Applikationsservern und versteckt zugleich
	deren komplexe Servlet\footnote{Eine Java-Klasse die HTTP Request's annimmt und mit einem entsprechenden Response beantworten
	kann.} \Abkuerzung{API} der Version 3.0. Ein besonderes Merkmal ist, dass das Framework je nach Applikationsserver und dem
	aufrufendem Browser die geeignete Transportmethode auswählen kann.
	\item [Spring]
	Das Spring Framework (kurz Spring) ist ein quelloffenes Framework für die Java-Plattform. Ziel des Spring Frameworks ist es, die
	Entwicklung mit Java/\Abkuerzung{Java EE} zu vereinfachen und gute Programmierpraktiken zu fördern. Spring bietet mit einem
	breiten Spektrum an Funktionalität eine ganzheitliche Lösung zur Entwicklung von Anwendungen und deren Geschäftslogiken; dabei
	steht die Entkopplung der Applikationskomponenten im Vordergrund. \citep{Spring}
	\item [jQuery]
	Bei jQuery handelt es sich um eine plattformunabhängige und umfangreiche Klassenbibliothek für die Scriptsprache JavaScript.
	Diese Open-Source-Software vereinfacht den Zugriff und die Manipulation des \Abkuerzung{HTML DOM} und wird
	insbesondere für ihre performante Programmierung wertgeschätzt. \citep[vgl.][]{jQuery}
	\item [LESS]
	Mit LESS wird \Abkuerzung{CSS} um dynamisches Verhalten erweitert. Damit kann die Wartbarkeit sowie Übersichtlichkeit 
	von \Abkuerzung{CSS} gesteigert werden. So sind \zB Vererbungen und die Definierung von Variablen mit unterschiedlichen
	Geltungsbereichen und Funktionen möglich. Es entspricht einer Annäherung zum objekt-orientierten Programmierparadigma.
\end{description}